\documentclass[11pt]{article}

\usepackage{fullpage}
\usepackage{amsmath, amssymb, bm, cite, epsfig, psfrag}
\usepackage{graphicx}
\usepackage{float}
\usepackage{amsthm}
\usepackage{amsfonts}
\usepackage{cite}
\usepackage{hyperref}
\usepackage{pgf,tikz}
\usepackage{enumitem}
\usepackage{mathtools}
\usepackage{siunitx}
\usepackage{../../styles/course}


\begin{document}

\title{Introduction to Hardware Design\\
Processor Interface:  Figures and Code Snippets}
\author{Profs. Sundeep Rangan, Siddharth Garg}
\date{}

\maketitle
\subsection{Vitis vs. HLS}

This is HLS code
\begin{systemverilogcode}
    module relu_lin (
      input logic int x,
      input logic int a,
      input logic int b,
      output logic int y
    );
      always_comb begin
        logic int mult_out, add_out;
        mult_out = x * a;
        add_out = mult_out + b;
        y = (add_out > 0) ? add_out : 0;
      end
    endmodule
  \end{systemverilogcode}

  This is equivalent C code:
  \begin{ccode}
    void relu_lin(int x, int a, int b, int &y) {
      int mult_out = x * a;
      int add_out = mult_out + b;
      if (add_out > 0) {
        y = add_out;
      } else {
        y = 0;
      }
    }
  \end{ccode}

  \subsection{Replacement code}
  \begin{ccode}
    // Compute x0 -> x1 -> x2
    easy_function1(x0, x1);  
    easy_function2(x1, x2); 

    // Compute x2 -> x3
    hard_function(x2, x3);  

    // Compute x3 -> x4
    easy_function3(x3, x4); 
  \end{ccode}

  \begin{systemverilogcode}
  module hard_function(...);
     ... 
  endmodule
  \end{systemverilogcode}

  \pagebreak
  \begin{ccode}
  void poly_stream(
      // Input and output data (AXI4 streams)
      hls::stream<float> &in_stream,
      hls::stream<float> &out_stream,

      // AXI4-Lite control
      float a0,
      float a1,
      float a2
  ) 
  \end{ccode}
  \end{document}

  

